\section{Life in the Noosphere}

\begin{quotex}
The fool does more harm by his ignorance than the wicked man by his wickedness.
\flright{\textsc{Ghazzali}}

To speak with spirits and angels was also common on this earth in ancient times, and for the same reason, namely, because they then thought much of heaven and little of the world. But that living communication with heaven in process of time was closed, as man from internal became external, or what is the same, as he began to think much about the world and little about heaven; and still more when he no longer believed that there is a heaven or a hell, nor that man in himself is a spirit which lives after death; for at this day it is believed that the body lives from itself, and not from its spirit; wherefore unless man now believed that he will rise again with the body, he would have no belief in the resurrection.
\flright{\textsc{Emmanuel Swedenborg}, \emph{Other Planets}}
\end{quotex}


You've heard of the collective unconscious, but what is more real is the collective conscious of humankind. Thought typically requires mediation through words --- in their many forms --- to propagate. However, it is obvious that certain thoughts propagate much more quickly than that; an idea is proposed and nearly immediately, millions of people become attached to that idea.

Beyond the human state, thoughts are shared directly without the intermediary of material forms. Swedenborg describes how that works in his encounters in the Jupiter state:

Presently angels of that planet came to me, and I was able to perceive from their speech with me, that they were altogether different from the angels of our earth; for their speech was not by words, but by ideas, which diffused themselves everywhere through my interiors.

Hence, ideas can be shared directly without anyone even being the wiser for it. These collective, shared ideas exist in what is called the Noosphere.

\paragraph{View from Kingston}


Several years ago, I visited a high ranking diplomat in Jamaica as her guest to the Marine Ball. She lived high on a hill in a posh neighbourhood with an armed sentry at the gate. Below, was the city of Kingston. I was kept awake at night from the sounds of the city, particularly the barking dogs and the howling hounds.

The thymos or incensive part of the soul in the soul can be a positive influence, although that same energy can provoke the soul to anger. In that case, they act like dogs, just barking through the night whether to protect property or to express fear and anxiety. The soul takes on the form of whatever rules it.

That what the human masses sound like from above, just incessant noise without purpose or any possibility or resolution.

\paragraph{View from Saturn}


The system of chakras is the microcosm that corresponds to the macrocosm of the planetary spheres. In particular, the Saturn sphere corresponds to the crown chakra at the top of the head, so it is the gateway between the individual state and the stars.

Hence. Saturn is associated both with melancholy and with mystical inspiration depending on the perspective: descent vs accent. Looking down toward earth, there is only confusion. People are unfree, not that they don't make choices, but rather that they are motivated by forces which are either unconscious or misunderstood. There are classes of people, differentiated by where their centre of consciousness resides.

\begin{enumerate}


\item \textbf{Lizards}: Their centre is the part of the soul guided by sensation. Their primary motivations are sex, fear, and hunger.

\item \textbf{Pigs}: Their centre is the appetitive part of the soul, where it becomes disordered and excessive.

\item \textbf{Dogs}: Their centre is the incensive part of the soul. They feel threatened and encounters lead to anxiety.

\item \textbf{Humans}. Their centre is in the intellectual part of the soul, so they are properly human. They use reason and are able to choose the good.
\end{enumerate}


From the Saturnine perspective, there seems to be no freedom in this scheme, so it seems predestined.

\paragraph{View from Earth}


After Adam and Eve were banished from Paradise, their perspective became earthbound. Their only source of knowledge derived from the senses and knowledge of higher worlds gradually become forgotten among humankind. At best, if a memory remained, it was the lowest form of memory, lacking any effective influence.

In this state, the Noosphere became closed to outside influences; The ``A'' influences include all those influences that originate in and are restricted to this closed system of thought. The assumption is that thought is empirically closed. In short, only empirical knowledge, and theories about it, count as valid. This leaves two open questions:

\begin{itemize}


\item \textbf{Consistency}: Is our empirical knowledge of the world consistent, or are there hidden contradictions.

\item \textbf{Completeness}: Is our knowledge of the world complete, at least in principle? Or are there things that can never be known. In other words, there are only ``A'' influences.
\end{itemize}


We can leave aside the question of completeness for now, although the quest for the ``theory of everything'' sees quite dubious. Consistency is a different matter. It would seem that at some point a consistent system should lead to certainty about certain things, since anything can be proved in an inconsistent system.

But that is precisely what we witness. On every point of significance to human life, there are multiple opinions continually at war with each other, with not conceivable way of resolving any differences.

The alternative is that the Noosphere is not closed. There are some, therefore, who are attuned to influences impinging upon the Noosphere from beyond. These influences, which we label ``B'' influences, are typically called ``Revelation''. Few still can hear the more subtle ``C'' influences, which are called ``esoteric'' in contrast to the exoteric ``B'' influences.

So the B and C influences show the way out, although not everyone is ready to accept them. The doubtful ones consider them just to be more ``A'' influences, although more absurd.

\paragraph{View from the Stars}


In the Divine Comedy, there is a long ladder from Saturn to the fixed stars, although it can be traversed in an instant. The call of Beatrice is all it takes. Listen for it.

What was understand in Saturn as \emph{Predestination} is now understood to be due to divine \emph{Providence}.

\paragraph{View from the Modern World}


In many regards the modern world is a blessing. To believe so requires trust in God and God's Providence. Despite its problems and degeneracy, there is more access to higher knowledge than every before. Books that used to be rare or difficult to find, can now be delivered right to your front door. In previous eras, you would be limited to knowing the people geographically close to you. Now you can meet people from across the globe who may be more spiritually attuned to you.

Medical advances lead to longer lives, which may provide needed time to understand esoteric texts, provided, of course, that time is not squandered. Almost nowhere in the world are people being executed for religious beliefs.

There is no cause for despair, since the world process follows cosmic laws and cycles. Moreover, the human state is merely one state of being, so it is possible to traverse the planetary spheres to reach the stars and heaven. Perhaps you may even find someone to make the journey with you!\footnote{The term ``noosphere'' should not be be assumed to refer to any previous use of that term.}

Dogs begin to barking and hounds begin to howl: \url{https://www.youtube.com/watch?v=mLrjhQ1-chw} 

\flrightit{Posted on 2021-06-06 by Cologero}

\begin{center}* * *\end{center}

\begin{footnotesize}
\begin{sffamily}
\texttt{Niko The Exile on 2021-06-23 at 16:42 said:}

Did you made a reference to Vernadski with the title or is this completely different thing from his views.

Mind you i only marginally know about Vernadski theories but i've heard about his term Noosphere.

\hfill

\texttt{Cologero on 2021-06-25 at 22:34 said:}

Niko, some readers of mysteries were in the habit of reading the end of the book first in order to discover ``whodunit''. Too bad I did not write a mystery, otherwise you would have noticed this disclaimer at the very end:

\begin{quotex}

The term ``noosphere'' should not be be assumed to refer to any previous use of that term.
\end{quotex}


But thanks for reading the title. That is something that 7 billion people will never bother to do.

\end{sffamily}
\end{footnotesize}

